\documentclass[11pt,a4paper]{article}
\usepackage[margin=2.4cm]{geometry}
\usepackage{amsmath,amssymb,amsthm,mathtools}
\usepackage{booktabs}
\usepackage{enumitem}
\usepackage{xcolor}
\usepackage{hyperref}
\usepackage{fontspec}
\usepackage{polyglossia}
\setmainlanguage{russian}
\setotherlanguage{english}
\setmainfont{Times New Roman}
\setmonofont{Consolas}[Scale=0.90]

\hypersetup{
  colorlinks=true,
  linkcolor=blue!70!black,
  citecolor=blue!70!black,
  urlcolor=blue!70!black,
  pdftitle={Ротаторное исчисление на кубооктаэдре, Часть IV: Геометрия планетарного дифференциала взаимно простых троек A + B = C, целочисленный вронскиан Мейсона-Стотерса-Безу, понижение степени на круге и межэтажный замок рода 10 в Lean 4},
  pdfauthor={Каратаев Дмитрий Александрович}
}

\theoremstyle{plain}
\newtheorem{theorem}{Теорема}[section]
\newtheorem{lemma}[theorem]{Лемма}
\newtheorem{proposition}[theorem]{Предложение}
\newtheorem{corollary}[theorem]{Следствие}
\theoremstyle{definition}
\newtheorem{definition}[theorem]{Определение}
\newtheorem{remark}[theorem]{Замечание}

\title{\textbf{Ротаторное исчисление на кубооктаэдре, Часть~IV:\\ Геометрия планетарного дифференциала взаимно простых троек $A + B = C$: целочисленный вронскиан Мейсона--Стотерса--Безу, понижение степени на круге и межэтажный замок рода~10 в Lean~4}}
\author{Д.\,А.~Каратаев}
\date{}

\begin{document}
\maketitle

\begin{abstract}
В четвёртой статье серии всякая взаимно простая аддитивная тройка $A + B = C$ ($A = a_1 a_2,\ B = b_1 b_2,\ C = c_1 c_2$, с этажным разложением $A = \widetilde{A}_k X_k^k,\ B = \widetilde{B}_k Y_k^k,\ C = \widetilde{C}_k Z_k^k$ и оценкой весов $\widetilde{A}_k\widetilde{B}_k\widetilde{C}_k \le \operatorname{rad}(ABC)^{k-1}$) формулируется как трёхвальный планетарный дифференциал с ортогональным $3\mathrm{D}$-репером $\mathcal{F}_3 = (\mathbf{r},\mathbf{\omega},\mathbf{r}\times\mathbf{\omega})\in\mathbb{Z}^{3\times 3}$ ($\mathbf{r} = (a_1,b_1,c_1)\perp\mathbf{\omega} = (a_2,b_2,-c_2)$), плоским законом скольжения шестерён Брахмагупты $C^2 + S_{AB}^2 = (A-B)^2 + M_{AB}^2 = (a_1^2+b_1^2)(a_2^2+b_2^2)$, невырожденной матрицей моментов Ганкеля--Вандермонда $-\det\mathcal{H}_3 = \widetilde{A}\widetilde{B}\widetilde{C}(X-Y)^2(Y-Z)^2(Z-X)^2 > 0$, проекторами Лагранжа--Вандермонда зазоров валов и унимодулярной $2\times 3$-матрицей партнёров Безу $\mathcal{B}_{2\times 3}$ (\texttt{thm8}, \texttt{thm9}).

Доказаны пять универсально $\forall$-квантифицированных диофантовых и геометрических теорем в кольце $\mathbb{Z}$:
(1)~обратная аннигиляция кубов Паде--Безу $Z^3 I_1 + A w^3(2Zv-Yw) \equiv \widetilde{B} I_0^3(Yw+Zv)$ (\texttt{thm12});
(2)~целочисленный аналог полиномиального вронскиана Мейсона--Стотерса на кубическом этаже $\widetilde{A}X^3+\widetilde{B}Y^3=\widetilde{C}Z^3$ через взаимную простоту Безу $\widetilde{A}Xv_0-\widetilde{B}Yu_0=1$, порождающий производную тройку $A^*+B^*=C^*$ ($A^*=\alpha X^2,\ B^*=\beta Y^2,\ C^*=\gamma Z^2$) с точным квадратным вронскианом валов $W^* = A B^* - B A^* \equiv X^2 Y^2 Z^2$ и определителем четвёртых степеней $-\det\mathcal{W}_{\mathrm{MS}} \equiv X^4 + Y^4 + Z^4$ (\texttt{thm13});
(3)~сквозная цепочка двустороннего схождения (\texttt{thm14}) с объёмно-радикальным тождеством $A B C = R_{\mathrm{full}} W^*$ для полного радикал-скелета с валами $R_{\mathrm{full}} = (\widetilde{A}\widetilde{B}\widetilde{C})(X Y Z)$ и безусловным понижением степени остатком на круге $0\le\alpha < \widetilde{A}X \implies (\widetilde{B}Y)C + (\widetilde{B}^2\widetilde{C})(Y^2 Z)\alpha = R_{\mathrm{full}}\beta \implies (\widetilde{B}Y)C \le R_{\mathrm{full}}\beta$, $C\le R_{\mathrm{full}}^2$ (при $\beta\le R_{\mathrm{full}}$) и $C\le R_{\mathrm{full}}^3$;
(4)~теорема о межэтажном сцеплении (\texttt{thm15}), показывающая, что в то время как один этаж $k=3$ имеет род $g(3)=1$, сцепление соседних этажей $k=2$ и $k=3$ на одних и тех же валах в силу $3-2=+1$ поднимает координаты $(U_6,V_6,W_6) = (X_2 Y_3 Z_3, Y_2 X_3 Z_3, Z_2 X_3 Y_3)$ на кривую 6-й степени Ферма--Сельмера $\widetilde{A}_2^3\widetilde{B}_3^2\widetilde{C}_3^2 U_6^6 + \widetilde{B}_2^3\widetilde{A}_3^2\widetilde{C}_3^2 V_6^6 = \widetilde{C}_2^3\widetilde{A}_3^2\widetilde{B}_3^2 W_6^6$ рода $g(6)=10\ge 2$;
(5)~универсальное стереографическое погружение Вейерштрасса (\texttt{thm16}) любой тройки $A+B=C\in\mathbb{Z}^3$ через $[0,512]/512 = [0,1]\cong\mathbb{R}$ в Квадрант~I $[0,128]$ ротатора $\mathbb{Z}_{512}$ с точным инвариантом четверти кольца $C\Delta_{\mathrm{rot}} - w I_{\mathrm{rot}} = 128$. Все теоремы верифицированы в ядре \textbf{Lean~4} (\texttt{0 sorryAx}, \texttt{0 Classical.choice}).

\medskip\noindent
\textbf{Ключевые слова:} аддитивные диофантовы тройки, $abc$-тройки, вронскиан Мейсона--Стотерса, аннигиляция Паде--Безу, кривые Ферма--Сельмера, стереографический гомеоморфизм Вейерштрасса, Lean~4.
\end{abstract}

\section{Репер планетарного дифференциала \texorpdfstring{$\mathcal{F}_3$}{F3}, матрица Ганкеля--Вандермонда \texorpdfstring{$\mathcal{H}_3$}{H3}, трёхканальная матрица Паде \texorpdfstring{$\mathcal{P}_3$}{P3} и перенос носителей \texorpdfstring{$M\rightleftarrows T$}{M <-> T}}
\label{sec:planetary-frame}

Следуя терминологии теории моделей А.\,И.~Мальцева (\emph{«Алгебраические системы»}~\cite{Malcev1970}, где \emph{носителем} / \emph{carrier} называется основное множество алгебраического сорта), мы рассматриваем систему как двусортную алгебраическую структуру $\langle\mathcal{M},\mathcal{T};\,\pi_{M\to T},\,\iota_{T\to M}\rangle$, связывающую \textbf{дискретный $M$-носитель} ($\mathbb{Z}_{512}$ с $L^1$-проекциями $\mathcal{C}_3,\mathcal{S}_3$) и \textbf{непрерывный евклидов $T$-носитель} ($L^2$-окружность $X^2+Y^2=R^2$ и континуальную прямую $\mathbb{R}$) через подъём Вейерштрасса $\pi_{M\to T}$, линейную ретракцию $\iota_{T\to M}$ и стереографический гомеоморфизм $[0,512]/512 = [0,1]\cong\mathbb{R}$.

\begin{theorem}[Ограниченное через радикал этажное погружение, окаймлённая $3\times 3$-матрица $\mathcal{M}_3$, ортогональный репер дифференциала $\mathcal{F}_3$, матрица Ганкеля--Вандермонда $\mathcal{H}_3$ и перенос носителей $M\rightleftarrows T$]
\label{thm:main-8}
Для всякой взаимно простой тройки $A + B = C$ ($\gcd(A,B,C)=1$) выполнены тождества (\texttt{thm8}, \texttt{thm9}):
\begin{enumerate}[label=\textbf{(\roman*)},leftmargin=*]
\item на каждом этаже $k\ge 2$ разложение $A = \widetilde{A}_k X_k^k,\ B = \widetilde{B}_k Y_k^k,\ C = \widetilde{C}_k Z_k^k$ удовлетворяет $\widetilde{A}_k\widetilde{B}_k\widetilde{C}_k \le \operatorname{rad}(ABC)^{k-1}$, а матрица $\mathcal{M}_3 = \begin{pmatrix} c_2 & a_2 & b_2 \\ a_1 & c_1 & 0 \\ b_1 & 0 & c_1 \end{pmatrix}$ имеет $\det\mathcal{M}_3 = c_1(C-A-B) = 0$ и союзную матрицу ранга~$1$ $(\mathcal{M}_3^{\mathrm{adj}})^2 = (c_1^2+C)\mathcal{M}_3^{\mathrm{adj}}$;
\item вектор радиусов шестерён $\mathbf{r} = (a_1,b_1,c_1)$ и вектор чисел оборотов $\mathbf{\omega} = (a_2,b_2,-c_2)$ ортогональны ($\mathbf{r}\cdot\mathbf{\omega} = 0$), $\det\mathcal{F}_3 = M_{BC}^2 + M_{CA}^2 + S_{AB}^2 \equiv \|\mathbf{r}\|_2^2\|\mathbf{\omega}\|_2^2 > 0$, $C^2 + S_{AB}^2 = (A-B)^2 + M_{AB}^2 \equiv (a_1^2+b_1^2)(a_2^2+b_2^2)$, а на этаже $k=3$ ($S_m = \widetilde{A}X^m + \widetilde{B}Y^m - \widetilde{C}Z^m$, $S_3 = 0$) матрица Ганкеля $\mathcal{H}_3$ имеет $-\det\mathcal{H}_3 \equiv \widetilde{A}\widetilde{B}\widetilde{C}(X-Y)^2(Y-Z)^2(Z-X)^2 > 0$ и проекторы Лагранжа--Вандермонда зазоров валов $\widetilde{A}X(X-Y)(X-Z) \equiv S_1 YZ - S_2(Y+Z)$, $\widetilde{B}Y(Y-X)(Y-Z) \equiv S_1 ZX - S_2(Z+X)$ и $\widetilde{C}Z(Z-X)(Z-Y) \equiv S_2(X+Y) - S_1 XY$;
\item трёхканальная кубическая матрица Паде $\mathcal{P}_3(A,B,C)$ имеет ядро $\mathcal{P}_3(A,B,C)(A,B,-C)^T = \mathbf{0}$ и след $\operatorname{Tr}\mathcal{P}_3(A,B,C) = A C(A+C)$, а по всем $512$ классам вычетов $x\in\mathbb{Z}_{512}$ степени $x^k\bmod 512$ сжимаются по лестнице $512\xrightarrow{k=4} 36\xrightarrow{k=8} 18\xrightarrow{k=16} 9\xrightarrow{k=32} 5\xrightarrow{k=64} 3\xrightarrow{k=128} 2 = \{0,1\}$ ($x^{128}\bmod 512 = x\bmod 2$).
\end{enumerate}
\end{theorem}

\section{Обратная аннигиляция Паде--Безу, целочисленный вронскиан Мейсона--Стотерса--Безу, понижение степени на круге, межэтажный замок рода~10 и стереографическое погружение в \texorpdfstring{$\mathbb{Z}_{512}$}{Z512}}
\label{sec:universal-theorems}

\begin{theorem}[Универсальные законы обратной аннигиляции Паде--Безу, вронскиана Мейсона--Стотерса, понижения степени на круге, межэтажного замка рода~10 и стереографического погружения в ротатор $\mathbb{Z}_{512}$]
\label{thm:main-12-16}
Для произвольных целых чисел в $\mathbb{Z}$ выполнены следующие универсально $\forall$-квантифицированные соотношения (\texttt{thm12}--\texttt{thm16}):
\begin{enumerate}[label=\textbf{(\roman*)},leftmargin=*]
\item \textbf{Обратная аннигиляция кубов Паде--Безу (\texttt{thm12}):}
  для $A = \widetilde{C}Z^3 - \widetilde{B}Y^3$, $I_0 = Yw - Zv$ и $I_1 = Yw(\widetilde{C}w^3+2\widetilde{B}v^3) - Zv(2\widetilde{C}w^3+\widetilde{B}v^3)$:
  \begin{equation}
    \label{eq:reverse-pade-bezout}
    Z^3 I_1 \;+\; A\,w^3(2Zv - Yw) \;\equiv\; \widetilde{B}\,I_0^3(Yw + Zv),\qquad
    Y^3 I_1 \;+\; A\,v^3(2Yw - Zv) \;\equiv\; \widetilde{C}\,I_0^3(Yw + Zv).
  \end{equation}
\item \textbf{Целочисленный кубический вронскиан Мейсона--Стотерса--Безу и определитель четвёртых степеней (\texttt{thm13}):}
  при $\widetilde{A}X^3 + \widetilde{B}Y^3 = \widetilde{C}Z^3$ и $\widetilde{A}Xv_0 - \widetilde{B}Yu_0 = 1$ подстановка $\alpha = u_0 Z^2 + t\widetilde{A}X$, $\beta = v_0 Z^2 + t\widetilde{B}Y$, $\gamma = u_0 X^2 + v_0 Y^2 + t\widetilde{C}Z$ даёт $A^*+B^*=C^*$ ($A^*=\alpha X^2,\ B^*=\beta Y^2,\ C^*=\gamma Z^2$) с минорами квадратов валов $(\widetilde{A}X\beta-\widetilde{B}Y\alpha,\ \widetilde{C}Z\alpha-\widetilde{A}X\gamma,\ \widetilde{C}Z\beta-\widetilde{B}Y\gamma) = (Z^2,Y^2,X^2)$, точным вронскианом
  \begin{equation}
    \label{eq:ms-wronskian}
    W^* \;:=\; A B^* - B A^* \;\equiv\; X^2 Y^2 Z^2,\qquad
    -\det\mathcal{W}_{\mathrm{MS}} \;\equiv\; X^4 + Y^4 + Z^4,
  \end{equation}
  и двусторонним схождением Крамера $w X^2 + \widetilde{B}\gamma = D_{\mathrm{clash}}Z$, $v X^2 + \widetilde{C}\beta = D_{\mathrm{clash}}Y$ ($D_{\mathrm{clash}} = \widetilde{C}w\beta - \widetilde{B}v\gamma$).
\item \textbf{Сквозная цепочка и безусловное понижение степени остатком на круге с полным радикал-скелетом валов (\texttt{thm14}):}
  для полного радикал-скелета с валами $R_{\mathrm{full}} := (\widetilde{A}\widetilde{B}\widetilde{C})(X Y Z)$ выполнено объёмное тождество $A B C \equiv R_{\mathrm{full}} W^*$ и точное круговое разложение
  \begin{equation}
    \label{eq:circle-remainder-rfull}
    (\widetilde{B}Y)\,C \;+\; (\widetilde{B}^2\widetilde{C})(Y^2 Z)\,\alpha \;\equiv\; R_{\mathrm{full}}\,\beta,\qquad
    (\widetilde{A}X)\,\beta \;=\; Z^2 + (\widetilde{B}Y)\,\alpha.
  \end{equation}
  Выбор неотрицательного остатка $0\le\alpha\le\widetilde{A}X-1$ на круге длины $\widetilde{A}X$ безусловно устраняет неопределённость знака: $(\widetilde{A}X)\beta \le Z^2 + (\widetilde{B}Y)(\widetilde{A}X-1)$ и $(\widetilde{B}Y)C \le R_{\mathrm{full}}\beta$, откуда $C \le R_{\mathrm{full}}^2 = \bigl((\widetilde{A}\widetilde{B}\widetilde{C})(X Y Z)\bigr)^2$ при $\beta\le R_{\mathrm{full}}$ и $C \le R_{\mathrm{full}}^3 = \bigl((\widetilde{A}\widetilde{B}\widetilde{C})(X Y Z)\bigr)^3$ для всех положительных валов и весов ($R_{\mathrm{full}}^3 - C \equiv C((\widetilde{A}\widetilde{B}\widetilde{C})^2 A B - 1)$).
\item \textbf{Межэтажное сцепление $k=2 \leftrightarrow k=3$ на кривой 6-й степени рода~10 (\texttt{thm15}):}
  сцепление этажа $k=2$ ($A = A_2 X_2^2,\ B = B_2 Y_2^2,\ C = C_2 Z_2^2$, род $g(2)=0$) и этажа $k=3$ ($A = A_3 X_3^3,\ B = B_3 Y_3^3,\ C = C_3 Z_3^3$, род $g(3)=1$) в силу $3-2=+1$ связывает каждый вал ($A A_3^2 X_3^6 = A_2^3 X_2^6$) и переводит $(U_6,V_6,W_6) = (X_2 Y_3 Z_3, Y_2 X_3 Z_3, Z_2 X_3 Y_3)$ на кривую 6-й степени рода $g(6) = \frac{5\times 4}{2} = 10 \ge 2$:
  \begin{equation}
    \label{eq:genus-10-curve}
    A_2^3 B_3^2 C_3^2\,U_6^6 \;+\; B_2^3 A_3^2 C_3^2\,V_6^6 \;=\; C_2^3 A_3^2 B_3^2\,W_6^6.
  \end{equation}
\item \textbf{Стереографическое погружение Вейерштрасса $\mathbb{Z}^3$ в Квадрант~I $[0,128]$ ротатора $\mathbb{Z}_{512}$ (\texttt{thm16}):}
  для любой тройки $A+B=C\in\mathbb{Z}^3$ и $A v - B u = 1$ ($u+v=w$) подъём Вейерштрасса $X_T = A^2-B^2,\ Y_T = 2AB,\ R_T = A^2+B^2$ удовлетворяет $X_T^2+Y_T^2=R_T^2$, $R_T+Y_T=C^2$, $(C^2+X_T,C^2-X_T)=(2CA,2CB)$ и $Y_T/(R_T+X_T) = B/A$, а в ячейке ротатора $\varphi_0\in[0,128]$ ($\mathcal{C}_3(\varphi_0)=128-\varphi_0,\ \mathcal{S}_3(\varphi_0)=\varphi_0,\ I_{\mathrm{rot}} = B\mathcal{C}_3(\varphi_0)-A\mathcal{S}_3(\varphi_0),\ \Delta_{\mathrm{rot}} = \mathcal{C}_3(\varphi_0)v-\mathcal{S}_3(\varphi_0)u$) выполняется точный инвариант четверти кольца:
  \begin{equation}
    \label{eq:quarter-ring-invariant}
    128 A \;=\; C\mathcal{C}_3(\varphi_0) - I_{\mathrm{rot}},\qquad
    128 B \;=\; C\mathcal{S}_3(\varphi_0) + I_{\mathrm{rot}},\qquad
    C\cdot\Delta_{\mathrm{rot}} \;-\; w\cdot I_{\mathrm{rot}} \;\equiv\; 128 \;=\; Q_\square(3).
  \end{equation}
\end{enumerate}
\end{theorem}

\section{Формальная верификация в Lean~4}
Формальное доказательство Теорем~\ref{thm:main-8} и~\ref{thm:main-12-16} (\texttt{thm8}, \texttt{thm9}, \texttt{thm12}, \texttt{thm13}, \texttt{thm14}, \texttt{thm15}, \texttt{thm16}) содержится в модуле \texttt{RotatorCalculusPaper.lean} в репозитории \url{https://github.com/AI-Guiders/scientific-papers} (\texttt{0 sorryAx}, \texttt{0 Classical.choice}).

\begin{thebibliography}{9}
\bibitem{Malcev1970}
А.\,И.~Мальцев, \emph{Алгебраические системы}, М.: Наука, 1970 (англ.\ пер.: A.\,I.~Mal'cev, \emph{Algebraic Systems}, Springer-Verlag, 1973).
\bibitem{Mason1984}
R.\,C.~Mason, \emph{Diophantine Equations over Function Fields}, Cambridge University Press, 1984.
\bibitem{DarmonGranville1995}
H.~Darmon and A.~Granville, ``On the equations $z^m = F(x,y)$ and $Ax^p + By^q = Cz^r$,'' \emph{Bull.~London Math.~Soc.} \textbf{27}(6) (1995), 513--543.
\bibitem{Moura2021}
L.~de~Moura and S.~Ullrich, ``The Lean~4 theorem prover and programming language,'' in \emph{CADE~28}, LNCS vol.~12699, Springer, 2021, pp.~625--635.
\end{thebibliography}

\end{document}
